\chapter{Homotopy Theory and the Fundamental Group}
\pagenumbering{arabic}

\begin{definition}
    If $f$ and $g$ are continuous maps from topological spaces $X$ and $Y$, we say that $f$ is \textui{homotopic} to $g$ if there exists a continuous map $F:X \times [0,1] \rightarrow Y$ such that:
        \begin{equation*}
        \begin{split}
            F(x,0) = f(x) \text{ and }F(x,1) = g(x)
        \end{split}
        \end{equation*}
    for all $x \in X$. We say $F$ is a \textui{homotopy} between $f$ and $g$, and write $f \cong g$. If $f$ is homotopic to a constant map, we say $f$ is \textui{nullhomotopic}.
\end{definition}

\begin{example}
    Let $f,g:X \rightarrow \bfR^n$ be continuous. Define:
        \begin{equation*}
        \begin{split}
            F(x,t) := (1-t)f(x) + tg(x). 
        \end{split}
        \end{equation*} 
    It is easy to see that $F$ is continuous, $F(x,0) = f(x)$, and $F(x,1) = g(x)$. The map $F$ is called the \textit{straight line homotopy}, because it moves the point $f(x)$ to the point $g(x)$ along a straight line.
\end{example}

\begin{example}
    Define:
        \begin{equation*}
        \begin{split}
            &i:S^1 \rightarrow \bfR^2 \setminus\{(0,0)\}, \h5 i(x) = x;\\
            &f:S^1 \rightarrow \bfR^2 \setminus\{(0,0)\},\h5 f(x) = (1,0).
        \end{split}
        \end{equation*}
    Then $i$ and $f$ are \textit{not} homotopic.
\end{example}

\begin{definition}
    A \textui{path} along a topological space $X$ is a continuous function $f:[0,1] \rightarrow X$. We call $f(0)$ the \textui{initial point} and $f(1)$ the \textui{terminal point}.
\end{definition}

\begin{definition}
    Two paths $f,g:[0,1] \rightarrow X$ are \textui{path homotopic} if their initial and terminal points are equal and there exists a continuous map $F:[0,1] \times [0,1] \rightarrow X$ such that:
        \begin{equation*}
        \begin{split}
            F(s,0) = f(s)\h9&\text{and} \h9 F(s,1) = g(s) \\
            F(0,t) = x_0 \h9\h8&\text{and} \h9 F(1,t) = x_1
        \end{split}
        \end{equation*}
\end{definition}